\chapter[1971 --- Sutherland et al.]{1971: Earl W. Sutherland, Jr.}
\label{ch:1971_earlwsutherlandjr}

\section*{Main Contribution}

\textit{Discovered the mechanism of hormone action---specifically, how epinephrine activates the enzyme adenyl cyclase to produce cyclic AMP as a second messenger.}

\section{Official Citation}

for his discoveries concerning the mechanisms of the action of hormones

\section{Background and Historical Context}

The 1971 Nobel Prize in Physiology or Medicine was awarded to Earl W. Sutherland, Jr. for his discoveries concerning the mechanisms of the action of hormones, especially adrenaline, through the intermediary formation of cyclic AMP. Sutherland's work represented a landmark in cell biology---the discovery of second messengers, a concept that fundamentally changed our understanding of how cells communicate with each other and respond to signals from their environment.

\subsection{The Hormone Signaling Puzzle}

By the mid-20th century, it was well established that hormones play important roles in regulating the functions of the body. Hormones are chemical messengers that are produced by endocrine glands and travel through the bloodstream to target cells, where they exert their effects. However, the mechanisms by which hormones exert their effects on target cells were still poorly understood.

The work of Sutherland would provide the first detailed understanding of how hormones signal target cells to produce their effects, through the discovery of cyclic AMP (cAMP) as a second messenger.

\subsection{Earl W. Sutherland, Jr.}

Earl W. Sutherland, Jr. was born on November 19, 1915, in Burlingame, Kansas. He studied at Washburn University, where he earned his bachelor's degree in 1937, and at Washington University School of Medicine in St. Louis, where he earned his medical degree in 1942. After completing his medical training, Sutherland worked at Washington University and later at Vanderbilt University, Case Western Reserve University, and the University of Miami, where he would spend most of his career.

Sutherland was a skilled biochemist and pharmacologist who was interested in the mechanisms of hormone action. His work on cyclic AMP would earn him the Nobel Prize.

\subsection{The Epinephrine Puzzle}

Sutherland's research began with a puzzling observation about the hormone epinephrine (adrenaline). When epinephrine is administered to animals, it causes an increase in blood sugar levels. This effect is due to the stimulation of glycogen breakdown (glycogenolysis) in the liver and muscles. However, the mechanisms by which epinephrine stimulates glycogenolysis were still unknown.

Sutherland hypothesized that epinephrine must act through some intermediary mechanism to stimulate glycogenolysis. This hypothesis led him to discover cyclic AMP, the first second messenger to be identified.

\section{The Discovery/Contribution}

Sutherland's work on hormone signaling involved the discovery of cyclic AMP as a second messenger and the elucidation of the mechanism by which hormones exert their effects on target cells.

\subsection{The Discovery of Cyclic AMP}

Sutherland's primary contribution was the discovery of cyclic AMP (cAMP), a small molecule that serves as a second messenger in many hormone signaling pathways.

\subsubsection{The Experimental Approach}

Sutherland used liver cells as his experimental system. He added epinephrine to liver cells and then measured the activity of glycogen phosphorylase, the enzyme that catalyzes the first step in glycogen breakdown. He discovered that epinephrine caused a rapid increase in the activity of glycogen phosphorylase, but only when the epinephrine was present outside the cell.

\subsubsection{The Breakthrough}

Sutherland discovered that when epinephrine binds to receptors on the surface of liver cells, it triggers the production of a small molecule that he called cyclic AMP (cAMP). This molecule then enters the cell and activates glycogen phosphorylase, leading to the breakdown of glycogen and the release of glucose into the bloodstream.

\subsubsection{The Second Messenger Concept}

Sutherland proposed that hormones act as first messengers, binding to receptors on the cell surface and triggering the production of second messengers (such as cAMP) inside the cell. The second messengers then activate intracellular signaling pathways that produce the cellular response.

This concept was significant because it explained how hormones could exert their effects on target cells without entering the cell themselves. It also explained how different hormones could produce different effects in different cell types, by activating different second messenger systems.

\subsection{The Adenylate Cyclase System}

Sutherland also discovered the enzyme adenylate cyclase, which catalyzes the conversion of ATP to cAMP. He showed that adenylate cyclase is located in the cell membrane and is activated by hormones binding to their receptors on the cell surface.

The adenylate cyclase system consists of three components:

\begin{itemize}
\item The hormone receptor: A protein in the cell membrane that recognizes and binds to the hormone
\item The G protein: A protein that transduces the signal from the receptor to the enzyme
\item Adenylate cyclase: The enzyme that catalyzes the production of cAMP
\end{itemize}

\subsection{Protein Kinase A}

Sutherland also discovered protein kinase A (PKA), an enzyme that is activated by cAMP and that phosphorylates other proteins, altering their activity. PKA is the primary effector of cAMP signaling and mediates many of the effects of cAMP in the cell.

\subsection{The Mechanism of Hormone Action}

Sutherland's work led to the elucidation of the complete mechanism of hormone action:

\begin{itemize}
\item The hormone (first messenger) binds to a receptor on the cell surface
\item The receptor activates a G protein, which in turn activates adenylate cyclase
\item Adenylate cyclase catalyzes the production of cAMP (second messenger) from ATP
\item cAMP activates protein kinase A, which phosphorylates target proteins
\item The phosphorylated proteins produce the cellular response
\end{itemize}

\section{Impact on Modern Medicine}

Sutherland's work on cyclic AMP and second messenger signaling has had a significant impact on modern medicine. His discoveries have contributed to the understanding and treatment of numerous diseases and have led to the development of many important drugs.

\subsection{Understanding Hormonal Disorders}

The understanding of second messenger signaling has been essential for understanding hormonal disorders, such as diabetes, thyroid disease, and adrenal disorders. These conditions often involve defects in hormone signaling pathways, and understanding these pathways has led to the development of new diagnostic tests and treatments.

\subsection{Drug Development}

The understanding of second messenger signaling has been essential for the development of many modern drugs. Numerous drugs work by affecting second messenger systems, including:

\begin{itemize}
\item Beta-blockers: Drugs that block beta-adrenergic receptors, reducing the effects of adrenaline and noradrenaline. These drugs are used to treat high blood pressure, heart failure, and other cardiovascular conditions.
\item Phosphodiesterase inhibitors: Drugs that inhibit the enzyme phosphodiesterase, which breaks down cAMP. These drugs increase the levels of cAMP in the cell and are used to treat conditions such as asthma and erectile dysfunction.
\item Adenylyl cyclase activators: Drugs that activate adenylate cyclase, increasing the production of cAMP. These drugs are used to treat conditions such as heart failure and pulmonary hypertension.
\end{itemize}

\subsection{Cancer Biology}

The understanding of second messenger signaling has also contributed to our understanding of cancer. Many cancer-causing genes encode components of second messenger signaling pathways, and mutations in these pathways can lead to uncontrolled cell growth. The principles established by Sutherland have been essential for understanding the mechanisms of cancer and for developing new treatments.

\subsection{Neuroscience}

The understanding of second messenger signaling has also contributed to neuroscience. cAMP plays important roles in the nervous system, including the regulation of neurotransmitter release, synaptic plasticity, and learning and memory. The principles established by Sutherland have been essential for understanding these processes and for developing new treatments for neurological disorders.

\subsection{Cardiovascular Medicine}

The understanding of second messenger signaling has also contributed to cardiovascular medicine. cAMP plays important roles in the regulation of heart rate and contractility, and drugs that affect cAMP signaling are widely used to treat cardiovascular conditions. The principles established by Sutherland have been essential for the development of these drugs.

\section{Legacy}

The legacy of Earl W. Sutherland, Jr. extends far beyond his Nobel Prize-winning work on cyclic AMP. His discoveries fundamentally changed our understanding of cell signaling and have contributed to numerous advances in medicine and biology.

\subsection{Foundation for Cell Signaling Research}

Sutherland's work established the foundations of modern cell signaling research. His discovery of cyclic AMP as a second messenger opened up an entirely new field of biology, and the principles he established continue to guide research in cell signaling today. The discovery of numerous other second messenger systems, including calcium, inositol trisphosphate, and diacylglycerol, was inspired by Sutherland's work.

\subsection{Other Second Messenger Systems}

Following Sutherland's discovery of cyclic AMP, researchers discovered numerous other second messenger systems that mediate the effects of hormones and other signaling molecules. These include:

\begin{itemize}
\item Calcium signaling: Calcium ions serve as second messengers in many cell types, regulating processes such as muscle contraction, neurotransmitter release, and gene expression
\item Inositol trisphosphate (IP3) and diacylglycerol (DAG): These lipid-derived second messengers are produced by the cleavage of phosphatidylinositol bisphosphate (PIP2) and regulate processes such as cell growth and differentiation
\item Cyclic GMP (cGMP): This nucleotide serves as a second messenger in many cell types, including photoreceptor cells in the eye and smooth muscle cells
\item Nitric oxide (NO): This gaseous signaling molecule activates guanylyl cyclase, leading to the production of cGMP
\end{itemize}

\subsection{G Protein-Coupled Receptors}

Sutherland's work on hormone receptors also contributed to the discovery and characterization of G protein-coupled receptors (GPCRs), a large family of cell surface receptors that mediate the effects of many hormones, neurotransmitters, and other signaling molecules. GPCRs are the targets of approximately 30 percent of all modern drugs, and understanding their structure and function has been essential for drug development.

\subsection{Advances in Medicine}

The understanding of second messenger signaling has contributed to numerous advances in medicine, including the development of new drugs for hormonal disorders, cancer, and cardiovascular disease. These advances have improved the lives of millions of patients and continue to drive progress in these fields.

\subsection{Inspiration for Future Researchers}

Sutherland's work continues to inspire researchers in cell signaling and pharmacology. His success in using elegant experiments to answer fundamental questions about hormone action demonstrates the power of basic research to improve human health. His example has motivated generations of scientists to pursue careers in cell signaling and pharmacology.

\subsection{Recognition and Honors}

In addition to the Nobel Prize, Sutherland received numerous other honors and awards for his work. He received the National Medal of Science in 1978 and was elected to the National Academy of Sciences. He also received honorary degrees from several universities and was recognized by the scientific community as one of the leading biochemists of his generation.

Sutherland's contributions to cell signaling and medicine are remembered and celebrated by researchers and clinicians around the world. His discovery of cyclic AMP and the second messenger concept continues to influence our understanding of cell signaling and has led to numerous advances in medicine and biology, ensuring that his legacy will endure for generations to come.


\section{Modern Developments}

Sutherland's second messenger work has led to modern signal transduction pharmacology. Understanding of cAMP signaling has enabled development of drugs targeting G protein-coupled receptors, which are targets of ~34% of approved drugs. Understanding of adenylyl cyclase has led to treatments for heart failure (milrinone) and pulmonary hypertension. Phosphodiesterase inhibitors (sildenafil) for erectile dysfunction and pulmonary hypertension build on his discoveries.
